
\documentclass{article}
\usepackage{colortbl}
\usepackage{makecell}
\usepackage{multirow}
\usepackage{supertabular}

\begin{document}

\newcounter{utterance}

\centering \large Interaction Transcript for game `cladder', experiment `full\_v1.5\_default', episode 4867 with qwen3.5{-}4b.
\vspace{24pt}

{ \footnotesize  \setcounter{utterance}{1}
\setlength{\tabcolsep}{0pt}
\begin{supertabular}{c@{$\;$}|p{.15\linewidth}@{}p{.15\linewidth}p{.15\linewidth}p{.15\linewidth}p{.15\linewidth}p{.15\linewidth}}
    \# & \multicolumn{2}{c}{Player} && \multicolumn{2}{c}{Game Master} \\
    \hline

    \theutterance \stepcounter{utterance}  
    & & & \multicolumn{4}{p{0.6\linewidth}}{
        \cellcolor[rgb]{0.9,0.9,0.9}{
            \makecell[{{p{\linewidth}}}]{
                \texttt{\tiny{[P1$\langle$GM]}}
                \texttt{You are an expert in causal inference. The following question is not a typical commonsense query, but rather a meticulously designed question created by a professor specializing in causal inference, intended to assess the students' mastery of the course content.} \\
\texttt{Imagine a self{-}contained, hypothetical world with only the following conditions, and without any unmentioned factors or causal relationships: Alice waking up has a direct effect on rainfall. Traffic has a direct effect on rainfall. The overall probability of waking up late is 9\%. The probability of waking up on time and high rainfall is 53\%. The probability of waking up late and high rainfall is 8\%. Is the chance of high rainfall smaller when observing waking up late?} \\
\texttt{Start your answer with 'Yes' or 'No', followed by additional reasoning or evidence to support your explanation.} \\
            }
        }
    }
     \\ \\

    \theutterance \stepcounter{utterance}  
    & \multicolumn{4}{p{0.6\linewidth}}{
        \cellcolor[rgb]{0.561,0.737,0.561}{
            \makecell[{{p{\linewidth}}}]{
                \texttt{\tiny{[P1$\rangle$GM]}}
                \texttt{Yes} \\
\\ 
\texttt{**Reasoning:**} \\
\\ 
\texttt{To determine whether the chance of high rainfall is smaller when observing waking up late, we must calculate the conditional probability of high rainfall given the state of waking up (late vs. on time).} \\
\\ 
\texttt{Let \$R\$ denote the event of high rainfall and \$W\$ denote the event of waking up late.} \\
\texttt{Let \$\textbackslash{}bar\{W\}\$ denote the event of waking up on time.} \\
\\ 
\texttt{We are given the following probabilities:} \\
\texttt{1.  \$P(W) = 0.09\$ (Probability of waking up late).} \\
\texttt{2.  \$P(\textbackslash{}bar\{W\} \textbackslash{}cap R) = 0.53\$ (Probability of waking up on time AND high rainfall).} \\
\texttt{3.  \$P(W \textbackslash{}cap R) = 0.08\$ (Probability of waking up late AND high rainfall).} \\
\\ 
\texttt{First, we calculate the total probability of high rainfall, \$P(R)\$, using the law of total probability:} \\
\texttt{\$\$P(R) = P(\textbackslash{}bar\{W\} \textbackslash{}cap R) + P(W \textbackslash{}cap R)\$\$} \\
\texttt{\$\$P(R) = 0.53 + 0.08 = 0.61\$\$} \\
\\ 
\texttt{Next, we calculate the conditional probability of high rainfall given that Alice woke up late, \$P(R|W)\$:} \\
\texttt{\$\$P(R|W) = \textbackslash{}frac\{P(W \textbackslash{}cap R)\}\{P(W)\} = \textbackslash{}frac\{0.08\}\{0.09\} \textbackslash{}approx 0.8889\$\$} \\
\\ 
\texttt{Then, we calculate the conditional probability of high rainfall given that Alice woke up on time, \$P(R|\textbackslash{}bar\{W\})\$. First, we need \$P(\textbackslash{}bar\{W\})\$:} \\
\texttt{\$\$P(\textbackslash{}bar\{W\}) = 1 {-} P(W) = 1 {-} 0.09 = 0.91\$\$} \\
\texttt{\$\$P(R|\textbackslash{}bar\{W\}) = \textbackslash{}frac\{P(\textbackslash{}bar\{W\} \textbackslash{}cap R)\}\{P(\textbackslash{}bar\{W\})\} = \textbackslash{}frac\{0.53\}\{0.91\} \textbackslash{}approx 0.5824\$\$} \\
\\ 
\texttt{**Comparison:**} \\
\texttt{*   \$P(R|W) \textbackslash{}approx 88.9\textbackslash{}\%\$} \\
\texttt{*   \$P(R|\textbackslash{}bar\{W\}) \textbackslash{}approx 58.2\textbackslash{}\%\$} \\
\\ 
\texttt{The probability of high rainfall is significantly **higher** (approximately 88.9\%) when observing Alice waking up late compared to when she wakes up on time (approximately 58.2\%).} \\
\\ 
\texttt{Therefore, the chance of high rainfall is **not** smaller when observing waking up late; it is larger. The answer to the specific question "Is the chance of high rainfall smaller when observing waking up late?" is **No**.} \\
            }
        }
    }
    & & \\ \\

    \theutterance \stepcounter{utterance}  
    & & & \multicolumn{2}{p{0.3\linewidth}}{
        \cellcolor[rgb]{0.9,0.9,0.9}{
            \makecell[{{p{\linewidth}}}]{
                \texttt{\tiny{[GM$|$GM]}}
                \texttt{Yes} \\
            }
        }
    }
    & & \\ \\

    \theutterance \stepcounter{utterance}  
    & & & \multicolumn{2}{p{0.3\linewidth}}{
        \cellcolor[rgb]{0.9,0.9,0.9}{
            \makecell[{{p{\linewidth}}}]{
                \texttt{\tiny{[GM$|$GM]}}
                \texttt{no} \\
            }
        }
    }
    & & \\ \\

    \theutterance \stepcounter{utterance}  
    & & & \multicolumn{2}{p{0.3\linewidth}}{
        \cellcolor[rgb]{0.9,0.9,0.9}{
            \makecell[{{p{\linewidth}}}]{
                \texttt{\tiny{[GM$|$GM]}}
                \texttt{game\_result = LOSE} \\
            }
        }
    }
    & & \\ \\

\end{supertabular}
}

\end{document}
